\section{What questions can we ask?}

Once we know that functions and sequences are related, we should also understand why they are not studied in exactly the same way. A sequence is a function, but its domain is usually \(\mathbb{N}\). A real function in calculus often has a domain that is an interval or a union of intervals. This changes the natural questions.

The purpose of this section is to organize those questions.

\subsection*{Questions about real functions}

For a real function
\[
f:D\to\mathbb{R},
\qquad
D\subseteq\mathbb{R},
\]
we often ask questions connected with its graph, values, and behavior as the input varies through real numbers.

Typical questions include:
\begin{itemize}
\item What is the domain?
\item What is the range?
\item What is \(f(a)\) for a given input \(a\)?
\item For which \(x\) do we have \(f(x)=0\)?
\item Where is the function positive or negative?
\item Is the function increasing or decreasing on an interval?
\item What happens to \(f(x)\) as \(x\) approaches a point?
\item Is the function continuous?
\end{itemize}

Not all of these questions are answered in this chapter. Some of them will appear later, especially limits and continuity. But the language needed to ask them begins here.

\subsection*{Questions about sequences}

For a sequence
\[
a:\mathbb{N}\to\mathbb{R},
\]
we ask questions adapted to natural-number inputs.

Typical questions include:
\begin{itemize}
\item What is the starting index?
\item What are the first few terms?
\item Is the sequence given by an explicit formula or recursively?
\item Is the sequence increasing or decreasing as \(n\) grows?
\item Is the sequence bounded above or below?
\item What happens to \(a_n\) as \(n\to\infty\)?
\item Does the sequence approach a limiting value?
\end{itemize}

The question about \(n\to\infty\) is especially important. Since the inputs of a sequence continue without end, we often want to know the long-term behavior of the terms.

\subsection*{A comparison}

\begin{center}
\begin{tabular}{lll}
\toprule
Object & Typical domain & Natural questions \\
\midrule
Real function & interval or subset of \(\mathbb{R}\) & graph, zeros, signs, continuity, limits as \(x\to a\) \\
Sequence & \(\mathbb{N}\) & terms, recursion, monotonicity, boundedness, limits as \(n\to\infty\) \\
\bottomrule
\end{tabular}
\end{center}

This table is not meant to separate functions and sequences completely. It is meant to show that different domains lead to different habits of investigation.

\subsection*{The common idea}

Despite these differences, the common idea remains the same. A function assigns outputs to inputs. A sequence also assigns outputs to inputs. The input set is special, but the structure is still functional.

For a real function, we may write
\[
f(x)
\]
for the value at input \(x\). For a sequence, we usually write
\[
a_n
\]
for the value at input \(n\). These notations look different, but they express the same idea: value at an input.

\begin{remarkbox}
The notation changes because the situation changes. But the underlying idea remains: one input gives one output.
\end{remarkbox}

\subsection*{Preparing for limits}

The next chapter will study limits and continuity. The present chapter prepares for that in two ways.

First, it gives the language of functions, values, domains, and graphs. Without that language, the expression
\[
\lim_{x\to a} f(x)
\]
has no clear meaning.

Second, it introduces sequences as functions on \(\mathbb{N}\). This prepares the expression
\[
\lim_{n\to\infty} a_n.
\]
Both limits involve the idea of approaching. But the input variable approaches in different ways: \(x\) may approach a real point through real values, while \(n\) goes through natural numbers toward infinity.

\begin{summarybox}
At the end of this chapter, remember:
\begin{itemize}
\item A function is an assignment from inputs to outputs.
\item The domain is part of the function.
\item The notation \(f(x)\) means the value of \(f\) at \(x\).
\item A graph is a picture of input-output pairs.
\item A sequence is a function whose domain is usually \(\mathbb{N}\).
\item Because sequences have discrete domains, we ask different questions about them.
\item The type of domain shapes the type of mathematics we do.
\end{itemize}
\end{summarybox}

Functions and sequences are not two unrelated topics placed side by side. Sequences are a special case of functions, and this viewpoint will help us understand limits, continuity, and later calculus.